但禁渔带来的恢复并非始终表现为生态健康改善。部分水域鱼类数量增长过快，水生植物和浮游动物被过度摄食，进而导致水体浑浊、底栖植被减少，部分水域甚至形成大面积"水下荒漠"。与此同时，长江江豚种群已回升至约1249头，2024年重新发现长江鲟自然繁殖产卵场；而鳄雀鳝、巴西龟等外来入侵物种已在鄱阳湖等流域支流存在稳定种群。十年禁渔确实取得了显著的资源恢复成效，恢复后的系统却未必更加健康，营养级失衡、江湖连通不足及入侵扩散依然存在。因此有必要在统一模型中分别计算"资源恢复"和"系统健康"。
\cite{report2022,report2023}。
